This paper addresses a practical challenge for long-form scientific writing with contemporary language models: producing section-level text that is coherent, auditable, and tightly coupled to explicit evidence.
Modern pretrained generative models can produce fluent outputs but may also generate unsupported or unverifiable assertions; the present work therefore adopts a conservative stance that prioritizes verifiability and traceability over unconstrained fluency.
Our central research question is: how does a graph-guided sectioning and prompting workflow affect narrative coherence and the fidelity of citation use during generation, relative to a linear prompting baseline?

To answer this question we propose a graph-aware prompting workflow that operationalizes two core ideas.
First, we represent a draft manuscript as a document graph in which nodes correspond to sections and subsections and edges encode scientific and rhetorical dependencies between them.
This representation supports a leaf-first scheduling strategy that writes dependent, evidence-rich lower-level nodes before synthesizing higher-level overview text, thereby reducing downstream incoherence and easing citation placement.
Second, we enforce machine-readable, evidence-constrained outputs for each section by (a) retrieving and ranking candidate evidence snippets scoped to the section, (b) injecting top-ranked provenance into prompts that require explicit citation markers for factual assertions, and (c) returning generated sections in a strict JSON schema that records text, structured citation entries, and provenance metadata.
At assembly time the JSON outputs are converted into a LaTeX manuscript with a refs.bib file and subjected to conservative validation checks—schema conformance, citation-coverage thresholds, and simple unsupported-claim heuristics—prior to human review and manual edits.

We emphasize conservative claims and auditability throughout the pipeline: every factual assertion in a generated section is required to reference a source in the retrieved evidence set, and the JSON outputs are intended to make post-hoc auditing, targeted edits, and reproducible assembly straightforward.
The JSON schema additionally records provenance identifiers and snippet anchors to support traceable edits and reviewer inspection, and it is designed to be human-interpretable as well as machine-parseable.
The paper contributes three practical artifacts for long-form generation workflows: (i) a graph-based section scheduler implementing leaf-first production, (ii) an evidence-injection protocol and ranking criteria used at section write time, and (iii) an audit-friendly JSON output format that integrates with standard LaTeX/bib toolchains and simple validation tools.
The remainder of the paper describes related work and positioning; details the graph construction, retrieval, prompt templates, and JSON schema; presents an experimental protocol comparing the graph-aware pipeline to a linear baseline on two topical domains; and discusses empirical observations, limitations, and future extensions.
We note as a caveat that, in the current draft, an exhaustive literature reconciliation was not completed and therefore claims of novelty are presented cautiously; a fuller related-work retrieval and citation set will be provided in subsequent revisions.